\section{Değerlendirme}
\label{CH:Degerlendirme}

\subsection{Başarı ölçütlerine göre durum}

Tez önerisindeki ölçütlerin ölçülen değerleri Tablo~\ref{tab:durum} içinde
özetlenmiştir.

\begin{table}[H]
\centering
\caption{Başarı ölçütleri ve ölçülen değerler}
\label{tab:durum}
\small
\begin{tabularx}{\textwidth}{@{}l l X l@{}}
\toprule
\textbf{Ölçüt} & \textbf{Hedef} & \textbf{Ölçülen} & \textbf{Durum} \\
\midrule
Top-1 doğruluk (AUTSL test) & $\geq$ \%85 & \%89,42 $\pm$ 0,52 (3 tohum); uygulamadaki int8 model \%89,82 & Sağlandı \\
Gerçek ortam doğruluğu & $\geq$ \%75 & \olculecek{gönüllü testi, Top-1} & -- \\
Model boyutu & $\leq$ 10 MB & 1,61 MB (int8 \texttt{.tflite}) & Sağlandı \\
İşleme hızı & $\geq$ 15 FPS & Yaklaşık 9 FPS (Redmi Note 10 Pro, önceki sürüm); yeni sürüm \olculecek{FPS} & Sağlanmadı \\
İnternet bağımlılığı & Tam çalışma & İnternet izni yok; uçak modunda \olculecek{deneme} & -- \\
\bottomrule
\end{tabularx}
\end{table}

\subsection{Hız}

Kare başına süre Redmi Note 10 Pro telefonda ön kamerayla, sınıflandırma modeli
olmadan ölçülmüştür (Tablo~\ref{tab:hiz}). Sınıflandırma modelinin kendisi bir işaret
bölütü için bir kez çalıştığından kare hızını belirleyen etken iki MediaPipe modelidir.

\begin{table}[H]
\centering
\caption{Yapılandırmaya göre kare hızı (Redmi Note 10 Pro, ön kamera)}
\label{tab:hiz}
\small
\begin{tabular}{@{}lrrrr@{}}
\toprule
\textbf{Yapılandırma} & \textbf{FPS} & \textbf{Kare başına (ms)} & \textbf{El (ms)} & \textbf{Gövde (ms)} \\
\midrule
Hata ayıklama sürümü, CPU, sıralı & $\sim$5 & 200--300 & 150--180 & 70--80 \\
Hepsi GPU, paralel & 7 & 120--130 & 90--120 & $\sim$100 \\
El GPU, gövde CPU, paralel & 9 & 110--120 & 70--90 & 80--100 \\
Gövde iki karede bir (yeni) & \multicolumn{4}{c}{\olculecek{tablette}} \\
\bottomrule
\end{tabular}
\end{table}

15 FPS hedefine ulaşılamamasının nedeni, her karede iki ayrı derin ağın (el ve gövde)
çalışmasıdır. Gövdenin iki karede bir çıkarılması bu yükü azaltmaktadır. Daha ileri
iyileştirmeler için analiz görüntüsünün küçültülmesi ya da el modelinin bir önceki
karedeki el konumundan başlatılması denenebilir. Öte yandan eğitimde uygulanan zamansal
seyreltme sayesinde model 10 FPS'lik girdide de çalışmaktadır: AUTSL test videoları
10 FPS'e seyreltildiğinde doğruluk yalnızca 0,8 puan düşmektedir (BiGRU, yüzde 88,00'den
yüzde 87,22'ye).

\subsection{Kayıtlı veri ile canlı kullanım arasındaki fark}

Proje boyunca yapılan canlı denemelerde tahmin güvenlerinin kayıtlı veriye göre
belirgin biçimde düştüğü gözlenmiştir. Bu farkı ayırmak için şu denemeler yapılmıştır:

\begin{itemize}
  \item \textbf{Öz-test:} Kayıtlı noktalar telefonda sınıflandırıldığında Python ile
        yüzde 99,5 aynı tahmin alınmıştır. Yani sorun ön işleme ya da modelin
        dönüşümünde değildir.
  \item \textbf{Video dosyasından telefon çıkarımı:} AUTSL videoları kamerasız olarak
        telefondaki MediaPipe zincirinden geçirildiğinde 14 videodan 10'u doğru
        tanınmıştır (yüzde 71; aynı kliplerde Python yüzde 81,7). Telefon ısındığı için
        deneme yarıda kesilmiş, örnek sayısı az kalmıştır.
  \item \textbf{Segmentleme:} 10 FPS'e seyreltilmiş AUTSL videolarında bölütleme
        kuralı doğruluğu bozmamaktadır (yüzde 89--90).
  \item \textbf{Kullanıcının işareti bilmesi:} İlk canlı denemeleri işaret dili
        bilmeyen ekip üyeleri yapmış, bazı denemelerde sözlükte olmayan işaretler
        (örneğin harfler) denenmiştir. Örnek hareket animasyonu bu nedenle eklenmiştir.
\end{itemize}

Bu bulgular farkın büyük ölçüde girdinin kendisinden (kişinin işareti doğru yapması,
kameraya uzaklık, ışık, kamera açısı) kaynaklandığını düşündürmektedir. Gönüllü
testinde toplanan kayıtlar \texttt{live\_analysis.py} ile AUTSL'ye göre karşılaştırılarak
bu etkenler sayısal olarak ayrıştırılacaktır. Sonuç: \olculecek{gönüllü testi analizi}.

\subsection{Gönüllü testi}

Gönüllü testi Ek~\ref{EK:Gonullu} içindeki protokole göre ekip dışından katılımcılarla
yapılmaktadır. Her katılımcı uygulamadaki test ekranında rastgele seçilmiş 30 kelimeyi
örnek hareketi izledikten sonra kameraya yapmaktadır.

\begin{table}[H]
\centering
\caption{Gönüllü testi sonuçları}
\label{tab:gonullu}
\begin{tabular}{@{}lrrr@{}}
\toprule
\textbf{Katılımcı} & \textbf{Deneme} & \textbf{Top-1 (\%)} & \textbf{Top-3 (\%)} \\
\midrule
\multicolumn{4}{c}{\olculecek{katılımcı sonuçları}} \\
\bottomrule
\end{tabular}
\end{table}
